\chapter{Deep Dive: Data Across Services --- Database-per-Service, Duplication, and the Outbox}
\label{ch:dd-consistency}

Chapter~7 established the rule: each service owns its database, and the
only road to another service's data is its API or its events. The rule
is easy to state and expensive to live with. This deep dive follows the
one piece of data the project copies between services --- a student's
name and e-mail --- from the moment it is written in auth-service's
Postgres to the moment it appears in course-service's, and asks what can
go wrong on the way. The answer is the dual-write problem, and the
repair is the transactional outbox.

Three questions organize the chapter:

\begin{enumerate}
  \item \textbf{What does database-per-service buy}, and what does it
  forbid --- joins, foreign keys, transactions --- that a monolith takes
  for granted?
  \item \textbf{How does duplicated data stay correct}, and what is the
  exact window in which the project is wrong today?
  \item \textbf{Can ``save the row and publish the event'' be made
  atomic} --- and if not, what is the closest thing to it?
\end{enumerate}

\section{Why this, and what it solves}

Open \code{docker-compose.yml}: two Postgres~17 containers,
\code{postgres-auth} holding \code{ts\_auth} and \code{postgres-course}
holding \code{ts\_course}, on separate volumes and ports. Nothing in
course-service knows the first one's address. The isolation buys four
things:

\begin{itemize}
  \item \textbf{Schema ownership.} Auth-service can add a
  \code{totp\_secret} column, rename \code{provider\_id}, or run a
  destructive Flyway migration without coordinating with anyone.
  Chapter~7's \code{V3} and \code{V4} migrations happened exactly that
  way.
  \item \textbf{Independent deploys.} A migration that locks
  \code{users} for a minute stalls logins, not course listings.
  \item \textbf{Independent scaling and tuning.} The course database
  will grow with submissions; the auth database will not. They can sit
  on different disks, get different connection pools, different backup
  schedules.
  \item \textbf{A hard stop on accidental coupling.} In a shared
  database the first \code{JOIN users ON student.user\_id = users.id}
  written under deadline pressure welds two teams' schemas together
  forever. Here the join is impossible, so it is never written.
\end{itemize}

It forbids exactly as much:

\begin{itemize}
  \item \textbf{No cross-service joins.} ``List students in course~7
  with their e-mails'' cannot be one query if e-mails live elsewhere.
  \item \textbf{No cross-service foreign keys.} \code{student.user\_id}
  in \code{V1\_\_create\_course\_tables.sql} is \code{BIGINT NOT NULL
  UNIQUE} and nothing more --- Postgres cannot check that user~42 exists
  in a database it cannot see.
  \item \textbf{No cross-service transactions.} A business operation
  that must change both databases cannot be rolled back as a unit.
\end{itemize}

\begin{decision}[copy the student, do not fetch the user]
Course-service needs a student's name and e-mail on every course roster
and in every assignment notification. Two designs were possible: call
auth-service each time (there is even a Feign client for it,
\code{AuthServiceClient.getUserById}, with a fallback that returns
\code{null}), or keep a local copy in the \code{student} table. The
project chose the copy, and the Feign client is dead code --- no class
in course-service calls it. The reasoning: a roster is read far more
often than a name changes; a synchronous call per row would make
course-service's availability a product of auth-service's (Chapter~6);
and the local copy can be indexed, joined, and paged with plain SQL. The
cost is that the copy can be \emph{stale}, and the rest of this chapter
is about how stale, and how to bound it. The decision flips when the
data changes often or must be authoritative at read time --- an account
balance, a permission --- where a call, or a cache with explicit
invalidation, is the honest design.
\end{decision}

\section{How a student reaches course-service}

Trace the real path. It starts in auth-service's \code{AuthService},
which is annotated \code{@Transactional} at class level:

\begin{lstlisting}[language=Java]
public void activate(String code) {
    User user = userRepository.findByActivationCode(code)
            .orElseThrow(() ->
                new IllegalArgumentException("Invalid activation code"));

    user.setActivated(true);
    user.setActivationCode(null);
    userRepository.save(user);                       (*@\co{1}@*)

    userActivatedPublisher.publish(user);            (*@\co{2}@*)
}
\end{lstlisting}

\code{UserActivatedPublisher} builds a \code{UserActivatedEvent} record
--- user id, e-mail, first and last name, role --- and calls
\code{kafkaTemplate.send("ts.user.activated", event.userId(), event)}.
On the other side, course-service's \code{UserActivatedListener}:

\begin{lstlisting}[language=Java]
@KafkaListener(topics = "ts.user.activated", groupId = "course-group")
public void onUserActivated(UserActivatedEvent event) {
    if (!ROLE_STUDENT.equalsIgnoreCase(event.role())) {
        return;
    }
    Long userId = Long.parseLong(event.userId());
    if (studentRepository.findByUserId(userId).isPresent()) { (*@\co{3}@*)
        return;
    }
    Student student = Student.builder()
            .userId(userId)
            .firstName(event.firstName())
            .lastName(event.lastName())
            .email(event.email())
            .build();
    studentRepository.save(student);                       (*@\co{4}@*)
}
\end{lstlisting}

\begin{description}
  \item[\co{1}] The row is dirty inside a transaction that has not yet
  committed. \code{save()} on a managed entity is a no-op until flush.
  \item[\co{2}] The event is handed to the Kafka producer \emph{inside}
  the same transaction, before the commit at method exit. Keep this line
  in mind; it is the whole problem.
  \item[\co{3}] The listener is idempotent by natural key: a second
  delivery of the same event finds the row and returns. This is why the
  project survives at-least-once delivery (Chapter~34) without a
  dedup table --- here.
  \item[\co{4}] The copy. From now on, course-service answers every
  roster query from this row and never asks auth-service again.
\end{description}

The consistency window is therefore: from the commit in
auth-service to the commit in course-service. On a quiet stack it is
milliseconds. With notification-service and course-service both in the
consumer path, and the broker on the same two-core box (Chapter~25), it
can be seconds. A teacher who enrols a student ``immediately'' after the
student activates may get \code{404 Student not found} from
\code{EnrollmentService.enroll} --- correct behaviour, surprising UX.

\subsection{What is not propagated at all}

Now look for the second event. There is none. \code{ts.user.activated}
fires once, at activation. A user who later changes their name or
e-mail in auth-service is never reflected in \code{student}; the
\code{DATABASE-DESIGN.md} note ``future: Kafka event propagates e-mail
changes'' is the acknowledgement. Assignment notifications will keep
going to the old address. This is the honest state of eventual
consistency in the project: \emph{eventual} for creation, \emph{never}
for updates. Section~\ref{sec:enhancing-outbox} adds the missing event.

\section{Pros and cons, candidly}

\begin{center}
\begin{tabular}{@{}p{0.3\linewidth}p{0.3\linewidth}p{0.3\linewidth}@{}}
\toprule
 & Shared database & Database-per-service (here) \\
\midrule
Cross-entity read & one \code{JOIN} & local copy, or two calls \\
Integrity & foreign keys & application code + events \\
Multi-entity write & one transaction & saga or outbox \\
Schema change & coordinate all owners & owner decides \\
Failure blast radius & one database, all services & one service \\
Staleness & none & seconds (creation), unbounded (updates) today \\
\bottomrule
\end{tabular}
\end{center}

The project's real weak points, named:

\begin{itemize}
  \item Update events do not exist; copies drift permanently.
  \item Save-and-publish is not atomic (next section).
  \item The Feign fallback returns \code{null}, so if anyone ever wires
  \code{AuthServiceClient} in, a circuit-open state will surface as a
  \code{NullPointerException} rather than a clear degraded response.
  \item No reporting path: ``how many activated students per course''
  needs data from both databases and there is no read model for it.
\end{itemize}

\section{The dual-write problem}
\label{sec:dual-write}

Line \co{2} above writes to two systems --- Postgres and Kafka --- that
share no transaction manager. Whatever order you put them in, one of
two failures is possible:

\begin{itemize}
  \item \textbf{Commit, then lose the event.} \code{send()} is
  asynchronous: it appends to a client-side buffer and returns a future.
  The method exits, the transaction commits, and the JVM is killed
  (\code{OOMKilled}, a rolling update, \code{kubectl delete pod}) before
  the buffer flushes. The user is activated; course-service never hears
  of it; the student cannot be enrolled, ever, and nothing logs an
  error.
  \item \textbf{Publish, then roll back.} The buffer flushes within
  milliseconds; the broker acks; then the commit fails --- a unique
  constraint, a connection dropped, a deadlock. Course-service creates a
  student for a user that does not exist. A \emph{phantom} row.
\end{itemize}

Both are rare per request and certain over a year. The interviewer's
version: ``you save an order and publish \code{OrderCreated}; how do you
guarantee both happen?'' The junior answer is ``put it in a
try/catch''. The senior answer starts with ``you cannot, with two
resources'', then names the outbox.

\begin{pitfall}[registration hangs for sixty seconds, then fails]
Symptom: with Kafka down, \code{POST /api/auth/register} does not fail
fast --- it hangs for a minute and returns 500. Cause: the producer must
fetch topic metadata before its first send; with no broker reachable it
blocks for \code{max.block.ms} (default 60\,000\,ms) and then throws.
The exception propagates out of \code{register()}, the transaction
rolls back, and the user is \emph{not} registered --- because the
mail-notification broker was down. A dependency that should have been
asynchronous has become a synchronous point of failure with a one-minute
timeout. The outbox below turns this into ``registered; e-mail delayed''.
\end{pitfall}

\begin{pitfall}[the copy that outlives the original]
Symptom: an assignment e-mail bounces from an address the user changed
months ago. Cause: \code{student.email} was written once by
\code{UserActivatedListener} and there is no \code{ts.user.updated}
event. Fix: emit one, and let the listener upsert. Section~\ref
{sec:enhancing-outbox} shows the emitter; the listener change is a
\code{findByUserId().map(update).orElseGet(create)}.
\end{pitfall}

\section{Enhancing the resolution: the transactional outbox}
\label{sec:enhancing-outbox}

The idea is to write the event into the \emph{same database}, in the
\emph{same transaction}, as the business row --- and let a separate
process copy it to Kafka afterwards. Postgres guarantees the row and the
event commit together or not at all; the relay guarantees the event
eventually reaches Kafka; the consumer's idempotency guarantees a
retried relay does no harm. Here is the full implementation for
auth-service.

The table, as Flyway migration \code{V5\_\_create\_outbox.sql}:

\begin{lstlisting}[language=SQL]
CREATE TABLE outbox (
    id             UUID PRIMARY KEY,
    aggregate_type VARCHAR(64)  NOT NULL,   -- 'user'
    aggregate_id   VARCHAR(64)  NOT NULL,   -- also the Kafka key
    event_type     VARCHAR(64)  NOT NULL,   -- 'UserActivated'
    topic          VARCHAR(128) NOT NULL,
    payload        JSONB        NOT NULL,
    created_at     TIMESTAMPTZ  NOT NULL DEFAULT now(),
    published_at   TIMESTAMPTZ,
    attempts       INT          NOT NULL DEFAULT 0
);
-- the relay's only query: oldest unpublished first
CREATE INDEX idx_outbox_unpublished
    ON outbox (created_at) WHERE published_at IS NULL;
\end{lstlisting}

The entity is deliberately dumb --- a row, not a domain object:

\begin{lstlisting}[language=Java]
@Entity @Table(name = "outbox")
@Getter @Setter @NoArgsConstructor
public class OutboxEvent {
    @Id private UUID id;
    private String aggregateType;
    private String aggregateId;
    private String eventType;
    private String topic;
    @JdbcTypeCode(SqlTypes.JSON)
    private String payload;                       // serialized event
    private Instant createdAt;
    private Instant publishedAt;
    private int attempts;
}

public interface OutboxRepository extends JpaRepository<OutboxEvent, UUID> {
    // SKIP LOCKED: two relay instances never pick the same row
    @Query(value = """
        SELECT * FROM outbox WHERE published_at IS NULL
        ORDER BY created_at LIMIT :n FOR UPDATE SKIP LOCKED""",
        nativeQuery = true)
    List<OutboxEvent> lockUnpublished(@Param("n") int n);
}
\end{lstlisting}

The publisher no longer touches Kafka. It writes a row, and because it
is called from inside \code{AuthService}'s transaction, the row commits
with the user:

\begin{lstlisting}[language=Java]
@Component
@RequiredArgsConstructor
public class UserActivatedPublisher {

    private final OutboxRepository outbox;
    private final ObjectMapper mapper;

    public void publish(User user) {
        UserActivatedEvent event = new UserActivatedEvent(
                String.valueOf(user.getId()), user.getEmail(),
                user.getFirstName(), user.getLastName(),
                user.getRole().name());
        OutboxEvent row = new OutboxEvent();
        row.setId(UUID.randomUUID());               (*@\co{1}@*)
        row.setAggregateType("user");
        row.setAggregateId(event.userId());          (*@\co{2}@*)
        row.setEventType("UserActivated");
        row.setTopic("ts.user.activated");
        row.setPayload(toJson(event));
        row.setCreatedAt(Instant.now());
        outbox.save(row);                            (*@\co{3}@*)
    }

    private String toJson(Object o) {
        try { return mapper.writeValueAsString(o); }
        catch (JsonProcessingException e) {
            throw new IllegalStateException(e);
        }
    }
}
\end{lstlisting}

The relay runs every second, takes a locked batch, sends each row
\emph{synchronously}, and marks it:

\begin{lstlisting}[language=Java]
@Component
@RequiredArgsConstructor
@Slf4j
public class OutboxRelay {

    private final OutboxRepository outbox;
    private final KafkaTemplate<String, String> kafka;

    @Scheduled(fixedDelay = 1000)
    @Transactional                                   (*@\co{4}@*)
    public void relay() {
        for (OutboxEvent row : outbox.lockUnpublished(100)) {
            ProducerRecord<String, String> rec = new ProducerRecord<>(
                    row.getTopic(), row.getAggregateId(), row.getPayload());
            rec.headers().add("eventId",
                    row.getId().toString().getBytes(UTF_8));    (*@\co{5}@*)
            rec.headers().add("eventType",
                    row.getEventType().getBytes(UTF_8));
            try {
                kafka.send(rec).get(10, TimeUnit.SECONDS);      (*@\co{6}@*)
                row.setPublishedAt(Instant.now());
            } catch (Exception e) {
                row.setAttempts(row.getAttempts() + 1);
                log.warn("outbox {} attempt {} failed: {}",
                        row.getId(), row.getAttempts(), e.toString());
                break;                                          (*@\co{7}@*)
            }
        }
    }
}
\end{lstlisting}

\begin{description}
  \item[\co{1}] The event id is minted \emph{here}, once, and survives
  every retry. It is the consumer's deduplication key (Chapter~34).
  \item[\co{2}] The aggregate id becomes the Kafka key, so all events
  about one user land on one partition, in order.
  \item[\co{3}] A plain \code{INSERT}, in the caller's transaction. If
  \code{activate()} rolls back, so does this row: no phantom events. If
  it commits, the row is durable before any Kafka involvement: no lost
  events.
  \item[\co{4}] The relay's transaction holds the \code{FOR UPDATE SKIP
  LOCKED} locks, so a second auth-service replica running the same
  scheduler takes different rows.
  \item[\co{5}] Headers carry identity and type without polluting the
  JSON payload the DTO classes deserialize.
  \item[\co{6}] The \emph{only} place in auth-service that waits on
  Kafka, and it is a background thread. HTTP requests never block on
  the broker again.
  \item[\co{7}] On failure stop the batch, keep ordering (row~2 must not
  be published before row~1 for the same key), and let the next tick
  retry. \code{attempts} is what you alert on.
\end{description}

Add \code{@EnableScheduling} to \code{AuthServiceApplication} and
change the \code{KafkaTemplate} to \code{<String, String>}, since the
payload is already JSON. The consumer sees exactly what it saw before.

Two rows of bookkeeping complete it: a nightly
\code{DELETE FROM outbox WHERE published\_at < now() - interval '7 days'}
keeps the table small, and the same \code{UserActivatedPublisher} shape
gives you \code{ts.user.updated} for free --- call it from wherever
profile changes are saved, with \code{eventType = "UserUpdated"}, and
have \code{UserActivatedListener} upsert instead of insert.

\begin{handson}[watch the outbox absorb a broker outage]
With the outbox in place, stop Kafka, activate a user, and look:
\begin{lstlisting}[language=bash]
$ docker compose stop kafka
$ curl -s -X POST localhost:8080/api/auth/activate?code=$CODE
{"message":"Account activated"}          # instant, no 60s hang
$ docker compose exec postgres-auth psql -U postgres ts_auth -c \
  "select event_type, attempts, published_at from outbox"
 event_type    | attempts | published_at
 UserActivated |        3 |
$ docker compose start kafka
# ...a few seconds later
 UserActivated |        3 | 2026-09-07 03:04:11.20+00
\end{lstlisting}
Then check course-service's log for \code{Provisioned student row}. The
activation was never at risk, and the student appeared as soon as the
broker did.
\end{handson}

\subsection{Change data capture: the industrial relay}

A polling relay is simple and adds up to one second of latency plus one
query per second per replica. At scale the relay is replaced by
\emph{change data capture}: Debezium tails Postgres's write-ahead log,
sees the \code{INSERT INTO outbox}, and publishes it to Kafka itself,
with no application scheduler and sub-100\,ms latency. The application
code is identical --- write the row --- which is why the outbox is worth
adopting early: the relay is the only part that changes.

\section{Sagas: when two services must both change}

Enrollment today is local --- \code{course}, \code{student}, and
\code{enrollment} are all in \code{ts\_course}, so
\code{EnrollmentService.enroll} is one ordinary transaction. The first
truly cross-service write the project will meet is \emph{account
deletion}: auth-service must remove the user; course-service must remove
the student, their enrollments, their submissions, and the submission
files in MinIO (Chapter~9). There is no transaction spanning that.

A \emph{saga} is a sequence of local transactions where each step
publishes an event that triggers the next, and each step has a
\emph{compensating} action if a later step fails:

\begin{enumerate}
  \item auth-service marks the user \code{DELETING} (not deleted) and
  writes \code{UserDeletionRequested} to the outbox.
  \item course-service deletes the student's data and files, then
  publishes \code{StudentDataDeleted} --- or, on failure,
  \code{StudentDataDeletionFailed}.
  \item auth-service, on the success event, hard-deletes the user; on
  the failure event, compensates by restoring \code{ACTIVE} and
  notifying an operator.
\end{enumerate}

This is \emph{choreography}: no coordinator, each service reacts to
events. It suits two or three steps. Past that --- deletion across
five services, with ordering constraints --- the event graph becomes
unreadable, and teams move to \emph{orchestration}: one component
(a saga orchestrator, or a workflow engine such as Temporal) holds the
state machine and sends \emph{commands} to each service. The trade is
one more service to run against a flow you can actually read.

Two properties every saga must have, and interviewers check for:
\emph{semantic locks} --- the \code{DELETING} state stops a user logging
in during the seconds the saga runs --- and \emph{idempotent steps}, so
the at-least-once delivery of Chapter~34 cannot delete twice or
compensate twice.

\section{Reads across services}

The remaining gap is reporting. ``Activated students per course, with
last login'' needs auth data and course data, and no service has both.
Three answers, in ascending cost:

\begin{itemize}
  \item \textbf{API composition.} The gateway or a small backend-for-
  frontend calls both services and joins in memory. Fine for one page,
  wrong for a nightly report over every course.
  \item \textbf{A read model.} A dedicated \code{reporting} database
  subscribes to \emph{every} domain event and maintains denormalized
  tables shaped for the queries. This is CQRS: writes go to the owning
  services, reads go to a model built from their events. The outbox
  above is the prerequisite --- without reliable events the read model
  drifts silently.
  \item \textbf{A warehouse.} CDC from every service database into an
  analytical store. The read model for people who are not you.
\end{itemize}

\section{What breaks at 3\,a.m.}

\begin{itemize}
  \item \textbf{postgres-course is down; users keep activating.}
  Without the outbox nothing is lost on the auth side, but
  \code{UserActivatedListener} throws on \code{save()}. Spring's default
  error handler retries ten times in quick succession, then \emph{logs
  and skips} the record (Chapter~34). The event is consumed, the
  student is never created, and the log line is the only trace. With
  the outbox the auth side is safe, but the consumer side still needs a
  retry topic --- the outbox fixes the producer, not the consumer.
  \item \textbf{An event is replayed.} Someone resets
  \code{course-group}'s offsets to recover from the outage above. Every
  \code{ts.user.activated} since retention start is redelivered.
  \code{findByUserId().isPresent()} makes each one a no-op. Idempotency
  by natural key is what makes replay a routine operation instead of an
  incident.
  \item \textbf{Ordering.} \code{UserActivated} and a future
  \code{UserUpdated} for the same user are keyed by \code{userId} and so
  share a partition; the update cannot overtake the activation. Events
  for \emph{different} users may interleave freely, which is fine ---
  there is no cross-user invariant. Change the key, or send from two
  different producers, and this guarantee is gone.
  \item \textbf{The relay stalls.} A poison row --- a payload the broker
  rejects for size --- makes \co{7} break on every tick, and rows behind
  it queue forever. Alert on \code{max(attempts)} and on
  \code{count(*) where published\_at is null}; after N attempts, move
  the row to a \code{outbox\_dead} table and let the queue drain.
\end{itemize}

\begin{decision}[when one database is the better call]
\code{docs/PLAN-A} describes the same product as a modular monolith:
one Spring Boot application, one \code{ts} database, and
\code{student.user\_id} as a real foreign key. Every problem in this
chapter --- the dual write, the drift, the saga, the read model ---
disappears there, replaced by a \code{JOIN} and a transaction. That is
not a small saving. The microservice layout (\code{PLAN-B}) earns its
cost only when the forbidding things are the point: teams that must
deploy without each other, data that must scale separately, or an
organization large enough that a shared schema \emph{would} be joined
across in anger. For a three-person team this project is deliberately
over-built, and Chapter~27's road ahead is honest about it. The senior
answer in an interview is not ``microservices'' or ``monolith''; it is
``a modular monolith until a boundary has a reason to be a network
boundary --- and when it does, the outbox from day one''.
\end{decision}

\section{Outside the project}

The outbox is the pattern behind Debezium's ``outbox event router'',
Axon's and Eventuate's transactional messaging, and the
\code{TransactionalEventPublisher} shapes in every serious event-driven
codebase. Sagas as orchestration are productized by Temporal, Camunda,
and AWS Step Functions; as choreography they are simply ``events with
compensations'' and a diagram someone must keep current. Netflix, Uber,
and every bank you can name run some combination of the three ---
outbox for reliability, sagas for multi-step writes, CQRS read models
for queries that cross ownership lines. The one thing none of them do is
what the project does today: call \code{send()} inside a transaction
and hope.

\section{Questions from real interviews}

\subsection*{Q: You need a report joining users and courses, and they live in two databases. How?}

First I ask how fresh and how big. For one screen, API composition: a
backend-for-frontend calls \code{/users?ids=} and \code{/courses}, joins
in memory, and that is fine up to a few hundred rows. For a real report
I build a read model: a \code{reporting} database that subscribes to
\code{ts.user.activated}, \code{ts.user.updated}, and the course events,
and maintains a denormalized \code{course\_roster} table shaped for the
query. That is CQRS --- the owning services keep their writes, the read
model answers cross-service reads. The prerequisite is reliable events,
which is why I put the outbox in first; a read model fed by best-effort
\code{send()} calls drifts and nobody notices until finance does. For
analytics at company scale the same idea is CDC into a warehouse. What I
refuse to do is give the reporting job a read replica of both databases
and a cross-database join --- that is the shared-database coupling the
architecture exists to prevent, back in through the side door.

\subsection*{Q: Two services must both succeed or neither. How, without two-phase commit?}

A saga. Each service does a local transaction and publishes an event;
the next step reacts; every step has a compensating action. In the
project, account deletion: auth-service sets \code{DELETING} and emits
\code{UserDeletionRequested}; course-service deletes data and emits
\code{StudentDataDeleted} or \code{...Failed}; auth-service either
hard-deletes or restores \code{ACTIVE}. Two things make it correct
rather than hopeful: a semantic lock --- the \code{DELETING} state
blocks logins for the seconds the saga runs --- and idempotent steps,
because Kafka will deliver at least once. I choose choreography for two
or three steps and orchestration (Temporal, Camunda, or a hand-written
state machine) once the event graph needs a diagram to follow. 2PC is
off the table not because it is wrong but because Kafka and Postgres do
not share a transaction coordinator, and holding locks across a network
call is how you get the 3\,a.m. page.

\subsection*{Q: An event was lost between the database commit and the Kafka send. How do you make that impossible?}

By never sending from the request thread. The transactional outbox: the
event is a row in an \code{outbox} table written in the same
transaction as the business row, so Postgres guarantees both or
neither. A relay --- a \code{@Scheduled} poller with \code{FOR UPDATE
SKIP LOCKED}, or Debezium tailing the WAL --- publishes it afterwards
and marks it. The relay is at-least-once, so the row carries an event id
that the consumer deduplicates on. In this project that also removes a
latent outage: today \code{kafkaTemplate.send()} inside
\code{AuthService.register} blocks for \code{max.block.ms} when the
broker is down, so a mail-broker outage makes registration fail after
sixty seconds. With the outbox, registration commits in milliseconds and
the e-mail is late instead of the signup being lost.

\subsection*{Q: How do you migrate a shared database into database-per-service safely?}

Strangler, not big bang. Step one: stop the bleeding --- find every
cross-schema join and foreign key and route them through the owning
service's API or a local copy, while the tables are still in one
database. That is the hard part and it can take months; the physical
split afterwards is a weekend. Step two: give each service its own
credentials and schema, so an accidental join fails at compile time.
Step three: dual-write or CDC the tables into the new database, verify
counts and checksums, flip reads, then writes, keep the old tables
read-only for a release, then drop. Throughout, the outbox goes in
\emph{before} the split, because the moment the join is gone you are
relying on events to keep copies correct. Chapter~7's ``two Postgres
servers, not two schemas'' is the end state; you reach it by
schema-per-service first.

\subsection*{Q: You duplicated Student data into course-service. How do you keep it correct?}

I accept that it will be stale and I bound the staleness. Creation
arrives on \code{ts.user.activated}; updates need a
\code{ts.user.updated} the project does not emit yet --- that is the
first fix. Both are keyed by \code{userId}, so they share a partition and
an update cannot overtake the creation. The listener upserts by natural
key, which makes replay safe. The events come from an outbox, so none
are lost; the consumer records processed event ids, so none are applied
twice. Then I measure: a nightly job compares a sample of
\code{student.email} against auth-service and alerts on drift, because
every one of the mechanisms above has a bug I have not found yet. What I
do not do is put the e-mail in the JWT --- that freezes it for the
token's lifetime and couples the schemas through the token contract
(Chapter~28).

\begin{quiz}
\begin{enumerate}
  \item Read \code{AuthService.activate}. Describe the two distinct
  failure orderings that leave auth-service and course-service
  disagreeing, and say which one produces a row that should not exist.
  \item Kafka is unreachable. Explain why \code{POST /register} today
  takes about a minute to fail and does not create the user, naming the
  producer setting responsible. What does the same request do after the
  outbox is in place?
  \item In the relay, why is \code{FOR UPDATE SKIP LOCKED} necessary,
  and what would go wrong with two auth-service replicas if it were
  replaced by a plain \code{SELECT}?
  \item Why does \code{break} (not \code{continue}) follow a failed
  send in the relay loop? Give the scenario the alternative corrupts.
  \item A student's e-mail changes in auth-service. State what
  course-service knows about it today, and list the three code changes
  that close the gap.
  \item Design account deletion across auth-service, course-service,
  and MinIO as a choreographed saga: name the states, the events, and
  the compensating action, and explain what the \code{DELETING} state
  protects against.
\end{enumerate}
\end{quiz}
